\documentclass[10pt,a4paper]{amsart}
\usepackage[margin=19mm]{geometry}
\usepackage{amsmath,amssymb,mathtools}
\usepackage[scr=rsfs,cal=euler]{mathalfa}
\usepackage{xfrac}
\usepackage[unicode,hidelinks,bookmarks=false]{hyperref}
\newcommand{\ie}{i.e.\ }
\newcommand{\cf}{cf.\ }
\newcommand{\resp}{respectively\ }
\newcommand{\Subsection}{\subsection}
\providecommand{\nobreakdash}{\nobreak\hbox{-}}
\setlength{\emergencystretch}{2em}
\multlinegap0pt
\usepackage{amssymb}
\usepackage{mathtools}
\usepackage[shortlabels]{enumitem}
\newtheorem{proposition}{Proposition}
\DeclareMathOperator{\Hess}{\mathrm{Hess}}
\DeclareMathOperator{\Ric}{\mathrm{Ric}}
\newcommand{\cA}{\mathcal{A}}
\newcommand{\cC}{\mathcal{C}}
\newcommand{\cL}{\mathcal{L}}
\title{Cylindrical Generalized Ricci Solitons in Three Dimensions}
\author{Miguel Pino Carmona}
\date{Source: arXiv:2606.03983v1 (CC BY 4.0)}
\begin{document}
\maketitle

\noindent\emph{Source and licence.} Cylindrical Generalized Ricci Solitons in Three Dimensions. Version arXiv:2606.03983v1 (CC BY 4.0), distributed by arXiv under the Creative Commons Attribution 4.0 International licence, http://creativecommons.org/licenses/by/4.0/, as recorded in the arXiv licence metadata for this exact version and shown on the abstract page https://arxiv.org/abs/2606.03983v1. Full licence notice: \texttt{licenses/arxiv-2606.03983-CC-BY-4.0.txt}.\par\smallskip

\noindent\textit{Editorial scope.} Counted unit: the article's proposition prop:dilatonconstant (source0.tex lines 652--654). The article's complete original proof of it is delivered with it (source0.tex lines 656--737, one \texttt{proof} environment of 82 lines). No mathematics is written, repaired or re-derived: the statement and the proof are the article's own lines, in the article's own notation, and no retained line is substituted or re-flowed.  Retained uncounted as the source-local context the proof needs: lines 85--88; lines 406--410; lines 427--429; lines 473--475; lines 494--522.  Nothing was omitted from any retained span, so the package carries no omission record.  No retained line contains a citation, so the wrapper carries no bibliography and no bibliography entry.  No retained line contains a figure environment, an image-inclusion command or drawing code, so no image is delivered and the package declares no asset dependency.  Editorial formatting only, in addition: an AMS \texttt{amsart} preamble that restates the article's own declarations and macros used by the retained lines, namely the environments \texttt{proposition} on separate counters, together with the standard packages those lines need.  The retained environments print in this document as Proposition 1 in order of appearance, whereas the article prints the same items with the numbering of its own full document.  The complete native source of the article as distributed in the arXiv e-print for this exact version is bundled unchanged as \texttt{sources/202016/source0.tex}.
\begin{align}
\Ric^g + \Hess^g f - \frac{1}{2}H \circ_g H &= 0 \,, \label{eq:einsteinGSS} \\
\Delta^g H + \cL_{\nabla f} H &= 0 \,. \label{eq:maxwellGSS}
\end{align}

\begin{align}
\phi'' &= -\frac{\psi'}{\psi} \phi' + f' \phi' - \frac{1}{2}k^2 e^{2f} \phi \,, \label{eq:phi''} \\
\psi'' &= -\frac{\phi'}{\phi} \psi' + f' \psi' - \frac{1}{2}k^2 e^{2f} \psi \,, \label{eq:psi''} \\
f'' &= \left(\frac{\phi'}{\phi} + \frac{\psi'}{\psi}\right) f' - 2\frac{\phi'\psi'}{\phi\psi} - \frac{1}{2}k^2 e^{2f} \,. \label{eq:f''}
\end{align}

\begin{equation}\label{eq:conservationlawreduced}
- f'' - \left( \frac{\phi'}{\phi} + \frac{\psi'}{\psi} \right) f' + (f')^2 - k^2 e^{2f} = \cC \,.
\end{equation}

\begin{equation}\label{eq:normalizedinitial}
\hat{\psi}(0)=1, \qquad \hat{f}(0)=0, \qquad \hat{k} = k \psi_0e^{f_0} > 0 \,.
\end{equation}

\begin{equation}\label{eq:boundaryconditions}
\phi(0)=0 \,, \quad \phi'(0)=1 \,, \quad \psi(0)=1 \,, \quad \psi'(0)=0 \,, \quad f(0)=0 \,, \quad f'(0)=0
\end{equation}

that ensure smoothness at the $t=0$ axis. This normalization will be used throughout the remainder of the paper.

However, once we fix $k$, not every solution satisfying the initial conditions \eqref{eq:boundaryconditions} is uniquely determined: since the system is singular at $t=0$, we rather have a family of solutions. We characterize this freedom using the reduced dilaton equation \eqref{eq:conservationlawreduced}, that is, the reduced conservation law associated to the system, through the constant $\cC$. Thus, for each fixed $k>0$, the initial conditions \eqref{eq:boundaryconditions} admit a one‑parameter family of smooth solutions, distinguished by the value of $\cC$; altogether we obtain a two‑parameter family of normalized solitons labeled by $k>0$ and $\cC$. As studied in Section \ref{sec:gradient generalized Ricci}, these solutions will also be complete when $\cC\ge0$.

%%%%%%%%%%%%%%%%%%%%%%%%%%%%%%%%%%%%%%%%%%%%%%%%%%%%%%%%%%%%%%%%%%%%%%%%%%%%
%%%%%%%%%%%%%%%%%%%%%%%%%%%%%%%%%%%%%%%%%%%%%%%%%%%%%%%%%%%%%%%%%%%%%%%%%%%%
%%%%%%%%%%%%%%%%%%%%%%%%%%%%%%%%%%%%%%%%%%%%%%%%%%%%%%%%%%%%%%%%%%%%%%%%%%%%
%%%%%%%%%%%%%%%%%%%%%%%%%%%%%%%%%%%%%%%%%%%%%%%%%%%%%%%%%%%%%%%%%%%%%%%%%%%%

\section{Explicit integration and global properties}
\label{sec:gradient generalized Ricci}

We now analyze the global structure of the gradient generalized Ricci solitons with diagonal cylindrical symmetry. Throughout this section, we work with the normalization \eqref{eq:normalizedinitial}.

%%%%%%%%%%%%%%%%%%%%%%%%%%%%%%%%%%%%%%%%%%%%%%%%%%%%%%%%%%%%%%%%%%%%%%%%%%%%
%%%%%%%%%%%%%%%%%%%%%%%%%%%%%%%%%%%%%%%%%%%%%%%%%%%%%%%%%%%%%%%%%%%%%%%%%%%%
%%%%%%%%%%%%%%%%%%%%%%%%%%%%%%%%%%%%%%%%%%%%%%%%%%%%%%%%%%%%%%%%%%%%%%%%%%%%
%%%%%%%%%%%%%%%%%%%%%%%%%%%%%%%%%%%%%%%%%%%%%%%%%%%%%%%%%%%%%%%%%%%%%%%%%%%%

\subsection{Calculation of explicit solutions}

For a diagonal cylindrical soliton $(g,H,f)$, we define the \emph{orbit area function}:
\begin{equation}\label{eq:areafunction}
\cA(t) = \sqrt{\det g(t)} = \phi(t)\psi(t) \,, \qquad \cA(0) = 0 \,, \qquad \dot{\cA}(0) = 1 \,.
\end{equation}

\begin{proposition}\label{prop:dilatonconstant}
Let $(\phi,\psi,f,k)$ be a gradient generalized Ricci soliton normalized as in \eqref{eq:normalizedinitial}. Then the solution has infinite diameter in the $t$-direction if and only if $\cC \ge 0$. In particular, for $\cC < 0$ the manifold has finite diameter in the $t$-direction and is incomplete.
\end{proposition}

\begin{proof}
From the Einstein equation $\eqref{eq:einsteinGSS}$ and the Maxwell equation $\eqref{eq:maxwellGSS}$, the metric components $\phi,\psi$ and the $H$-flux satisfy the reduced Einstein equation \eqref{eq:phi''}–\eqref{eq:f''}. As shown in Section \ref{sec:gradient generalized Ricci}, \eqref{eq:phi''} and \eqref{eq:psi''}, determine $\phi$ and $\psi$ uniquely in terms of the variable $\tau(t) = \int_0^t e^{f(s)} \,ds$, independently of the dilaton equation. Explicitly,
\begin{equation*}
\phi(\tau) = \frac{2}{k}\sin(k\tau/2),\qquad \psi(\tau) = \cos(k\tau/2),
\end{equation*}
and consequently $\cA(\tau) = \phi\psi = \frac{1}{k}\sin(k\tau)$. The function $\cA(\tau)$ is positive precisely for $0 < k\tau < \pi$, so $\tau$ can never exceed $\pi/k$; the maximal possible interval is $[0,\pi/k)$. The actual domain of $\tau$ will be $[0,\tau_{\max})$, where $\tau_{\max} \le \pi/k$ is determined by the condition $y(\tau) = e^{f(\tau)} > 0$.

Recall that $\cA'/\cA=\phi'/\phi+\psi'/\psi$. Then, using $d/d\tau=e^{-f}\,d/dt$, we compute
\begin{equation*}
f' = e^f\dot{f} \,, \qquad f'' = e^{2f}(\ddot f + \dot{f}^2) \,, \qquad \frac{\cA'}{\cA} = e^f \,\frac{\dot{\cA}}{\cA} = e^f k\cot(k\tau) \,.
\end{equation*}

Thus, we can write \eqref{eq:f''} as
\begin{equation*}
0 = f'' - \left(\frac{\phi'}{\phi} + \frac{\psi'}{\psi}\right) f' + 2\frac{\phi'\psi'}{\phi\psi} + \frac{1}{2}k^2 e^{2f} = e^{2f}(\ddot{f} + (\dot{f})^2) - e^{2f} k\cot(k\tau) \,\dot{f} \,,
\end{equation*}
hence
\begin{equation*}
0 = e^f(\ddot{f} + (\dot{f})^2) - e^f k\cot(k\tau) \,\dot{f} = \frac{d^2}{d\tau^2}(e^f) + k\cot(k\tau) \,\frac{d}{d\tau}(e^f) \,.
\end{equation*}

Introducing the variable $y=e^f$, we obtain the linear second-order ODE
\begin{equation*}
\ddot{y} - k\cot(k\tau)\,\dot{y} = 0 \,,
\end{equation*}
with initial conditions $y(0)=1,\ \dot{y}(0)=0$. The general solution is
\begin{equation}\label{eq:ysolution}
y(\tau) = A + D\cos(k\tau) \,, \qquad A+D=1 \,.
\end{equation}

Now, expressing \eqref{eq:conservationlawreduced} in terms of $\tau$ yields
\begin{equation*}
- e^{2f}\left(\ddot{f} + k\cot(k\tau)\dot{f} + k^2\right) = \cC \,,
\end{equation*}
and, using $y=e^f$ and
\begin{equation*}
f = \ln y \,, \qquad \dot{f} = \frac{\dot{y}}{y} \,, \qquad \ddot{f} = \frac{\ddot{y}}{y} - \frac{(\dot{y})^2}{y^2} \,,
\end{equation*}
we can write it as
\begin{equation*}
y\ddot{y} - (\dot{y})^2 + k\cot(k\tau)y\dot{y} + k^2y^2 = - \cC \,.
\end{equation*}

Since $\ddot{y} - k\cot(k\tau)\,\dot{y} = 0$, we can further simplify the equation to
\begin{equation*}
2k\cot(k\tau)y\dot{y} - (\dot{y})^2 + k^2y^2 = - \cC \,.
\end{equation*}

Finally, expanding using \eqref{eq:ysolution} and simplifying, we find
\begin{equation*}
k^2(A^2 - D^2) = - \cC \,,
\end{equation*}
and, using $A+D=1$, we simplify it to
\begin{equation}\label{eq:Dcondition}
k^2(1 - 2D) = - \cC \,.
\end{equation}

Hence
\begin{equation*}
D = \frac12\left(1 + \frac{\cC}{k^2}\right) \,, \qquad A = \frac12\left(1 - \frac{\cC}{k^2}\right) \,,
\end{equation*}
and therefore
\begin{equation}\label{eq:ysolutionfinal}
y(\tau) = \frac12\left(1-\frac{\cC}{k^2}\right) + \frac12\left(1+\frac{\cC}{k^2}\right)\cos(k\tau) \,.
\end{equation}

We now determine the completeness of the metric from the behavior of $t(\tau) = \int_0^\tau \frac{ds}{y(s)}$. If $\cC < 0$, then $y(\tau)=\tfrac12\left(1+\frac{|\cC|}{k^2}\right) + \tfrac12\left(1-\frac{|\cC|}{k^2}\right)\cos(k\tau)$. If $|\cC|\le k^2$ we have $y(\tau)\ge \frac{|\cC|}{k^2}>0$; if $|\cC|>k^2$ we have $y(\tau)\ge1$. In either case $y$ is bounded away from zero on $[0,\pi/k]$. Consequently the integral $t(\tau) = \int_0^\tau \frac{ds}{y(s)}$ converges to a finite limit $t_{\max}$ as $\tau\to\pi/k$. Thus the geodesic distance $t$ is bounded; the manifold has finite diameter and is incomplete.

If $\cC = 0$, then $A = D = 1/2$, and $y(\tau) = \frac12(1+\cos(k\tau)) = \cos^2(k\tau/2)$. Here $\tau_{\max} = \pi/k$, and
\begin{equation*}
t(\tau) = \int_0^\tau \frac{ds}{\cos^2(ks/2)} = \frac{2}{k}\tan(k\tau/2) \longrightarrow \infty \quad \text{as } \tau \to \pi/k \,,
\end{equation*}
so the manifold is complete.

Finally, if $\cC > 0$, then $D > A$ and $y(0)=1>0$, while $y(\pi/k) = A - D = -\cC/k^2 < 0$. By continuity there exists a unique $\tau_0 \in (0,\pi/k)$ such that $y(\tau_0)=0$. Since $D \neq 0$, this zero is simple. As $\tau \to \tau_0^-$, $1/y(\tau)$ has a simple pole, hence
\begin{equation*}
t(\tau) = \int_0^\tau \frac{ds}{y(s)} \longrightarrow \infty \quad \text{as } \tau \to \tau_0^- \,.
\end{equation*}
Therefore the geodesic distance $t$ attains all values in $[0,\infty)$, and the manifold is geodesically complete. The domain of $\tau$ is $[0,\tau_0)$, with $\tau_0 = \frac{1}{k}\arccos(-A/D) < \pi/k$.

Combining the three cases, we conclude that the solution is complete precisely when $\cC \ge 0$, and incomplete for $\cC < 0$. This finishes the proof.
\end{proof}

\end{document}
